\section{Design decisions}
\label{sec:design-decisions}
Key architectural decisions were made based on the requirements and constraints identified in the analysis phase.

\subsection{Block Editor: Ggoogle Blockly}
\begin{table}[h]
  \centering
  \label{tab:block-editor-selection}
  \caption{Block Editor Selection}
  \begin{tabularx}{0.9\textwidth}{|l|X|X|l|}
    \hline
    Option                & Pros                                                 & Cons                                     & Decision   \\\hline
    Google Blockly        & Mature, documented, accessible, customizable toolbox & Larger bundle size                       & \checkmark \\\hline
    Custom implementation & Full control, smaller size                           & Development time, maintenance burden     &            \\\hline
    Scratch Blocks        & Familiar Scratch aesthetics                          & Less flexible, tied to Scratch ecosystem &            \\\hline
  \end{tabularx}
\end{table}

Rationale: Blockly provides proven reliability, accessibility features and the flexibility to define per-exercise toolboxes.
The bundle size trade-off is acceptable given lazy loading optimizations.

\subsection{Code Execution: Client-Side Sandbox}
\begin{table}[h]
  \centering
  \label{tab:code-execution}
  \caption{Code Execution Options}
  \begin{tabularx}{0.9\textwidth}{|l|X|X|l|}
    \hline
    Option                       & Pros                               & Cons                                                                & Decision   \\\hline
    Client-side (iframe sandbox) & No server load, immediate response & Limited isolation compared to server-side, potential security risks & \checkmark \\\hline
    Server-side (VM/Container)   & Better isolation                   & Added complexity, latency                                           &            \\\hline
    WebAssembly sandbox          & Best isolation                     & Significant development effort                                      &            \\\hline
  \end{tabularx}
\end{table}

Rationale: For educational code from children, client-side execution with loop traps and API whitelisting provides sufficient security.
Server-side execution would add complexity and latency without proportional benefits for this use case.

\subsection{Database: SQLite with Drizzle ORM}
\begin{table}[h]
  \centering
  \label{tab:database-selection}
  \caption{Database Selection}
  \begin{tabularx}{0.9\textwidth}{|l|X|X|l|}
    \hline
    Option      & Pros                         & Cons                          & Decision   \\\hline
    SQLite      & No setup, embedded, portable & Single-writer limitation      & \checkmark \\\hline
    PostgresSQL & Scalable, concurrent writes  & Requires server, more complex &            \\\hline
  \end{tabularx}
\end{table}

Rationale: QLite offers the best balance for single-school deployments. The Drizzle ORM
provides type-safe queries and migration management. The schema is designed to be portable,
enabling future migration to PostgreSQL if needed.


\section{Architectural Overview}
\label{sec:design-overview}
The system follows a layered architecture separating presentation, business logic and data access concerns.

\begin{figure}[h]
  \centering
  \includegraphics[width=\textwidth]{architecture.excalidraw.png}
  \caption{System Architecture Overview}
  \label{fig:architecture-overview}
\end{figure}

\subsection{Component Descriptions}
\subsubsection*{Frontend Layer}
The frontend is built with SvelteKit and comprises three main interfaces:
\begin{itemize}
  \item \textbf{Student Interface:} Displazs exercises, the Blockly workspace, canvas visualization and grading feedback.
        Designed for simplicity and age-appropriateness.
  \item \textbf{Teacher Interface:} Enables exercis creation and editing, course organization and progress monitoring. Includes visual editors for paths and grids.
  \item \textbf{Admin Interface:} Manages users, system configuration and plugin administration.
\end{itemize}

\subsubsection*{Exercise Engines}
Exercise engines encapsulate the state and behavior for each exercise type:
\begin{itemize}
  \item \textbf{Turtle:} Manages position, angle, pen state, and path history. Commands include \textit{move(distance)} \textit{turn(degrees)}, and pen control.
  \item \textbf{GridRobot:} Manages grid position, direction and inventory. Command sinclude directional movement, collection and interaction,.
  \item \textbf{I/O:} Manages input consumption and output accumulation for stdin/stdout-style exercises.
\end{itemize}

\subsubsection*{Sandbox}
The sandbox provides secure code execution through:
\begin{itemize}
  \item An ifram with \textit{sandbox="allow-scripts"} attribute
  \item PostMessage-based communication protocol
  \item Whitelisted API surface exposing only exercise-specific functions
  \item Loop trap injection detecting infinite loops
  \item Timeout enforcement terminating long-running executions
\end{itemize}

\subsubsection*{Graders}
Type-specific grades evaluate student solutions.
\begin{itemize}
  \item \textbf{Turtle Grader:} Compares final state (position, angle) or command sequencees against expected values with configurable tolerances.
  \item \textbf{I/O Grader:} Compares program output against expected output with nromalization options.
\end{itemize}


\subsection{Exercise Type Framework}
The exercise framework is dsigned for extensibility through a common base class and grading interface.
\subsubsection*{Abstract Base: Canvas2D}
All visual exercies types inherit from Canvas2D, which provides:
\begin{itemize}
  \item Reactive state for position coordinates
  \item Command logging for grading
  \item Reset and comparison methods
  \item Canvas rendering hooks
\end{itemize}

\subsubsection*{Concrete Implementations}
\textbf{Turle} Extends Canvas2D with angle, pen state, and move/turn methods.
\textbf{GridRobot} Extends Canvas@d with discrete grid positioning, direction and inverntory management.
\textbf{I/O} Implements a non-visual exercise type with input consumption and output generation capabilities.

New exercise types can be added by:
\begin{enumerate}
  \item Implementing a class extending Canvas2D
  \item Defining corresponding Blockly blocks and generators
  \item Implementing a grade for the new type
\end{enumerate}

\section{Data Model}
\label{sec:design-datamodel}
The data model or database schema supports exercies, courses, user progress and audit logging.
\subsection{Core Tables}
\begin{itemize}
  \item \textbf{users} Stores user credentials, roles (student, teacher, author, admin ) and preferences.
  \item \textbf{courses} Contains course metadata, ordered exercies list via \textbf{courseExercises} junction table, and enrollment via \textbf{courseUsers}.
  \item \textbf{exercises} Stores exercise definitions including type (io/turlte/robot), JSON content (localized title/description), JSON config (type-specific grader settings), toolbox XML, starter XML and hints.
  \item \textbf{attempts} Records student attempts with workspace XML, grading resutls, score and timestamp.
  \item \textbf{auditLogs} Tracks all admin actions like login, account cration, password reset for accountability and auditing.
  \item \textbf{exerciseVersions} Stores snapshosts of exercises for versioning support
  \item \textbf{translations} Custom i18n table: locale/kez/value rows used for runtime language lookup
\end{itemize}

\subsection{Relationships}
\begin{itemize}
  \item Courses contain ordered lists of exercies (many-to-many with ordering)
  \item Users have roles determining access to CMS and admin features
  \item Attempts link users to exercises with temporal tracking
\end{itemize}

\section{Security Model}
\subsection{Threat Model}
Primary threats addressed by the security model include:
\begin{itemize}
  \item Malicious student code attempting DOM access or infinite loops
  \item Cross-site scripting through user-generated content
  \item Unauthorized access to exercises or administrative functions
\end{itemize}

\subsection{Mitigitations}
\begin{itemize}
  \item \textbf{Iframe Sandbox} The \textit{sandox="allow-scripts"} attribute prevents navigation, from submission and plugin access.
  \item \textbf{Loop Trap} Command counting detects and terminates infinite loops.
  \item \textbf{API Whitelist} Only exercise-specific functions are exposed; no access to window, document, or network APIs.
  \item \textbf{CSP Headers} Strict Content Security Policy prevents inline scripts and external resource loading.
  \item \textbf{Input Validation} Zod schema validation on all API inputs prevents injection.
  \item \textbf{Role-Based Access} Teachers cannot access admin functions; students cannot access CMS.
\end{itemize}

